\section{Implementation and Diagnosis}
\label{sec:implementation}

Reproducing \mcpilot{} \cite{turcato2025mcpilot} from the \mcpilco{}
codebase \cite{mcpilco} revealed four failure modes that are individually
sufficient to prevent any learning from scratch.
None of these defects cause an obvious runtime crash in the general case; each
manifests instead as degenerate-but-silent optimiser behaviour.
We document each below using the structure: \textbf{Symptom $\to$ Root cause $\to$
Fix $\to$ Evidence}.

% ─────────────────────────────────────────────────────────────────────────────
\subsection{Broken Policy Gradient Path}
\label{sec:fix_gradient}

\textbf{Symptom.}
During early training runs the optimiser raised
\texttt{RuntimeError: element 0 of tensors does not require grad and does not have a
grad\_fn} on the very first backward pass.
Even after patching the immediate exception, cost values were constant across
optimisation steps, confirming that no gradient reached the policy parameters.

\textbf{Root cause.}
The \mcpilco{} codebase separates particle propagation into two stages:
\texttt{apply\_policy()} computes the commanded speed $u$ from the RBF policy and
passes it to the simulator, and \texttt{data\_to\_gp\_input()} pre-processes the
augmented state before GP inference.
\texttt{data\_to\_gp\_input()} strips all input dimensions that do not appear in the
GP's training index set — including the policy-derived speed channel.
As a result, the computation graph from policy weights $\boldsymbol{\theta}$ through
$u$ to the GP output was severed; \texttt{cost.backward()} found no path to
$\boldsymbol{\theta}$.

\textbf{Fix.}
The commanded speed $u$ is embedded \emph{directly} into the particle initial state
rather than being passed as a separate GP input.
Given target bearing $\phi = \operatorname{atan2}(P_y - y_0,\; P_x - x_0)$ and fixed
elevation $\alpha = 35^\circ$, the 3D release velocity is:
\begin{align}
  v_x &= u\cos\alpha\cos\phi, \label{eq:vx} \\
  v_y &= u\cos\alpha\sin\phi, \label{eq:vy} \\
  v_z &= u\sin\alpha.         \label{eq:vz}
\end{align}
All three operations (\texttt{atan2}, \texttt{cos}, \texttt{sin}) have well-defined
PyTorch autograd implementations.
The velocity vector $\mathbf{v} = [v_x, v_y, v_z]^\top$ is written into particle
state dimensions 3:6 before the rollout begins, so every GP call and cost evaluation
has an unbroken computation graph back to $\boldsymbol{\theta}$.

\textbf{Evidence.}
After this fix, \texttt{cost.grad\_fn} is non-\texttt{None} at every optimisation
step, and policy parameters update from the first backward pass onward.
The cost trajectory exhibits monotone descent from step~1, confirming gradient flow.

% ─────────────────────────────────────────────────────────────────────────────
\subsection{Post-Landing GP Hallucination}
\label{sec:fix_landing}

\textbf{Symptom.}
Before any fix, the per-trial cost plateaued at $\approx 0.998$ regardless of target
distance or policy speed, and optimisation produced no improvement over the
$N_{\mathrm{exp}} = 5$ exploration throws.
Inspecting particle states revealed that some particles had $z \ll 0$
(up to $-3$\,m) at the terminal timestep $T = 1.0$\,s.

\textbf{Root cause.}
The simulation horizon in the paper is $T = 1.0$\,s with $T_s = 0.01$\,s
($= 100$ timesteps).
A tennis ball thrown at the maximum speed $u_M = 3.5$\,m/s lands in approximately
$0.58$\,s under realistic drag.
The \mcpilco{} codebase continued propagating all particles through the remaining
$\approx 42$ timesteps after landing, feeding underground positions
($z < 0$, $v_z$ still negative or zero) as GP inputs.
These states lie far outside the GP's training support, so the posterior mean is
governed by the zero prior and the posterior variance is maximal.
The resulting pseudo-dynamics drove particles to arbitrary positions, and the terminal
cost $c_T = 1 - \exp(-\|\mathbf{e}_T\|^2 / \ell_c^2)$ evaluated these garbage states,
yielding $c_T \approx 1$ for all particles regardless of $\boldsymbol{\theta}$.

\textbf{Fix.}
A Boolean \emph{landed} mask $\mathbf{m} \in \{0,1\}^M$ is maintained across
timesteps.
Once a particle's $z$-coordinate crosses zero (detected at each step), its state is
frozen: $\mathbf{s}^{(i)}_{t+1} \leftarrow \mathbf{s}^{(i)}_t$ for all subsequent $t$.
The landing instant is refined by linear interpolation between the last airborne
and first underground states, giving the exact ground-crossing position used for cost
evaluation.
This matches the procedure described in \S4.2.2 of \cite{turcato2025mcpilot} but
was absent from the \mcpilco{} codebase.

\textbf{Evidence.}
With the landed mask active, terminal particle $z$-coordinates are all $\leq 0$ and
the GP is never queried on underground inputs.
Baseline cost drops from $\approx 0.998$ to $\approx 0.88$ after the exploration
phase, providing a non-trivial gradient signal from trial~1.

% ─────────────────────────────────────────────────────────────────────────────
\subsection{Cost Lengthscale Bootstrapping Failure}
\label{sec:fix_lc}

\textbf{Symptom.}
After fixing the gradient path and landing mask, training costs oscillated near $0.88$
for all $N_{\mathrm{opt}} = 1500$ optimisation steps of every trial, with
per-step cost changes below $10^{-4}$.
Policy parameters moved, but landing errors showed no improvement between trials.

\textbf{Root cause.}
The paper specifies a cost lengthscale $\ell_c = 0.1$\,m in the squared-exponential
terminal cost:
\begin{equation}
  c(\mathbf{s}_T) = 1 - \exp\!\left(-\frac{\|\mathbf{e}\|^2}{\ell_c^2}\right),
  \label{eq:cost}
\end{equation}
where $\mathbf{e} = [x_T - P_x,\; y_T - P_y]^\top$ is the landing error.
With $\ell_c = 0.1$\,m, $c \geq 0.99$ whenever $\|\mathbf{e}\| \geq 0.22$\,m, and
the gradient $\partial c / \partial \|\mathbf{e}\| = (2\|\mathbf{e}\|/\ell_c^2)\,
\exp(-\|\mathbf{e}\|^2/\ell_c^2)$ is smaller than $10^{-3}$ for
$\|\mathbf{e}\| > 0.3$\,m.
Exploration throws against randomly-placed targets produced landing errors in the
range $0.3$--$1.5$\,m, placing all particles firmly in the saturated,
near-zero-gradient region of~\eqref{eq:cost}.
The GP therefore learned correct dynamics from exploration data, but the cost
function could not distinguish any two bad policies, so the optimiser stalled.

The necessary condition for a non-zero bootstrap gradient is:
\begin{equation}
  \ell_c \;\gtrsim\; \operatorname{median}\!\left(\{\|\mathbf{e}^{(k)}\|\}_{k=1}^{N_{\mathrm{exp}}}\right).
  \label{eq:lc_rule}
\end{equation}

\textbf{Fix.}
$\ell_c$ is set to $0.5$\,m.
At this value $c$ varies meaningfully over the interval $[0,\;1]$\,m, covering the
full range of exploration errors, and $|\partial c / \partial \|\mathbf{e}\||$ is
non-negligible up to $\approx 1$\,m.

\textbf{Evidence.}
With $\ell_c = 0.5$\,m, the cost function exhibits a gradient magnitude above
$10^{-2}$ for all exploration-phase errors observed.
Optimisation converges to hit rate $5/10$ (baseline seed) within
$N_{\mathrm{opt}} = 1500$ steps per trial.
Reverting to $\ell_c = 0.1$\,m while keeping all other fixes in place restores the
stalled-cost behaviour, confirming this as an independent failure mode.

% ─────────────────────────────────────────────────────────────────────────────
\subsection{Simulation Horizon and Indexing Bugs}
\label{sec:fix_horizon}

Two further bugs related to simulation horizon and data-structure indexing caused
silent logical errors with the original paper parameters.

\textbf{Symptom A: T\_control mis-specification.}
The MC-PILCO interface expects a \emph{number-of-steps} integer for the control
horizon.
The paper reports $T = 1.0$\,s.
Passing $T = 1.0$ directly caused the internal computation
\texttt{control\_horizon = int(T / Ts) = int(1.0 / 0.01) = 100}... to be applied
to what the codebase treated as already-an-integer, yielding an effective horizon of
$T \cdot (1/T_s)^2 = 10{,}000$ steps.
The resulting rollout allocated a $10{,}000 \times M$ particle tensor, exhausting RAM
and crashing on any machine with less than \needdata{measured peak RAM} of available
memory.

\textbf{Fix A.}
The argument is passed as \texttt{T\_control = int(T / Ts) = 50} (steps, for
$T = 1.0$\,s and $T_s = 0.02$\,s), matching the codebase's expected type.

\textbf{Symptom B: opt\_steps\_list index out of range.}
\mcpilco{} indexes a pre-allocated list \texttt{opt\_steps\_list} by
\texttt{trial\_index}, where \texttt{trial\_index} starts at
$N_{\mathrm{exp}} - 1 = 4$ and increments.
The list was initialised with $N_{\mathrm{opt}}$ entries (one per optimisation trial),
so \texttt{trial\_index = 4} immediately accessed an out-of-range position, raising
\texttt{IndexError} at the start of the first learning trial.

\textbf{Fix B.}
\texttt{opt\_steps\_list} is allocated with length
$N_{\mathrm{exp}} + N_{\mathrm{trials}}$, ensuring valid access for all indices.

\textbf{Horizon optimisation.}
Beyond the above correctness fixes, the simulation horizon was reduced from
$T = 1.0$\,s ($T_s = 0.01$\,s, 100 steps) to $T = 0.7$\,s ($T_s = 0.02$\,s,
35 steps).
The maximum ball flight time at $u_M = 3.5$\,m/s is $\approx 0.58$\,s; no particle
can be airborne beyond $0.7$\,s.
Combined with the landed-freeze mask (Section~\ref{sec:fix_landing}), this halves the
number of GP evaluations per rollout and yields an approximately $3\times$ wall-clock
speedup with no accuracy regression, verified by comparing landing-error distributions
across seeds.

\textbf{Remark on exploration policy.}
Algorithm~1 of \cite{turcato2025mcpilot} initialises exploration with an
analytically-derived policy (Eq.~13 therein) that computes near-optimal parabolic
throws for each target.
Applied in our setting this produced all exploration throws concentrated in a narrow
speed band $[3.1,\;3.4]$\,m/s, causing the GP to overfit to a small region of
velocity space.
Hit rate \emph{dropped} from 50\% (random exploration) to 20\% (paper's analytical
exploration) when both were tested under identical seeds.
We therefore use uniform-random exploration across $[u_m, u_M] = [0.75, 3.5]$\,m/s
for all experiments, and document the analytical policy as a failed experiment.
